\section{Related work}
\label{related_work}
In this section, we review the current research on perception in outdoor environments and explore offline RL methods applied to navigation tasks. Finally, we discuss the existing methods for covert navigation planning.
\subsection{Perception for Outdoor Navigation}
Accurate perception of the environment is crucial for autonomous navigation in complex outdoor scenes. Traditional approaches rely on supervised learning techniques, such as semantic segmentation \cite{sathyamoorthy2022terrapn} and object detection \cite{sathyamoorthy2020frozone}, to identify traversable terrain and potential cover. However, these methods require extensive manual annotation and may not generalize well to new environments or robot platforms \cite{papadakis2013terrain, weerakoon2022terp}.
Recent works have explored self-supervised learning \cite{sathyamoorthy2022terrapn} and adversarial training \cite{fan2020learning} to reduce the reliance on human annotation. While these techniques have shown promise, they still struggle to fully capture the complex properties and dynamics of real-world environments.

% In contrast to these approaches, our framework leverages multi-modal sensory information, specifically LiDAR intensity and height data, to generate high-fidelity cover maps and potential threat maps. This allows the robot to accurately perceive and reason about the environment's properties and potential cover without relying on extensive manual annotation.

In contrast to these approaches, our framework leverages the rich information provided by LiDAR point cloud data to generate high-fidelity height maps, cover maps, and threat maps. By processing the LiDAR data, our framework enables the robot to accurately perceive and reason about the environment's properties, including terrain elevation, potential cover locations, and threat positions, without relying on extensive manual annotation or simplified representations. 

% This LiDAR-based perception approach allows EnCoMP to capture detailed environmental information and facilitate intelligent decision-making for covert navigation in complex outdoor environments.

\subsection{Offline Reinforcement Learning for Navigation}
Offline reinforcement learning is a promising approach for learning navigation policies from previously collected datasets, mitigating the challenges associated with online learning in real-world environments \cite{levine2020offline, fujimoto2019off}. Recent works have demonstrated the effectiveness of offline reinforcement learning for navigation tasks in simulated environments \cite{mandlekar2020iris, hansen2021generalization}. Hansen et al. \cite{hansen2021generalization} introduced a generalization through offline reinforcement learning framework that learns a navigation policy from a diverse dataset of trajectories in simulated environments. Shah et al.  presented ReViND ~\cite{shah2023offline}, the first offline RL system for robotic navigation that can leverage previously collected data to optimize user-specified reward functions in the real-world. The system is evaluated for off-road navigation without any additional data collection or fine-tuning, and it is shown that it can navigate to distant goals using only offline training from this dataset. Another work by Li et al. ~\cite{hao2023avoidance} proposes an efficient offline training strategy to speed up the inefficient random explorations in the early stage of navigation learning.

Our approach builds upon these advancements by introducing an offline reinforcement learning framework specifically designed for covert navigation in complex outdoor environments. By learning from a diverse dataset collected from real-world settings, our approach can effectively capture the complex dynamics and uncertainties present in these environments and generate robust navigation strategies.

\subsection{Covert Navigation}
Covert navigation is a critical capability for robots operating in hostile or sensitive environments, where the primary objective is to reach a designated goal while minimizing the risk of detection by potential threats \cite{ajai2019survey, niroui2019deep}. Existing approaches to stealth navigation often rely on heuristic methods \cite{tews2004avoiding} or potential field-based techniques \cite{marzouqideceptive} to identify potential cover and minimize exposure. However, these methods may not effectively capture the complex dynamics and uncertainties present in real-world environments. Recent works have explored the use of reinforcement learning for stealth navigation ~\cite{mendoncca2018reinforcement}. While these approaches have shown promise in learning adaptive strategies, they often rely on simplified environmental representations and do not fully exploit the rich sensory information available in real-world settings.

Our framework addresses these limitations by integrating LiDAR-based perception and offline reinforcement learning, enabling robots to navigate covertly and efficiently in complex outdoor environments. By leveraging the rich environmental information captured by LiDAR sensors and training a robust policy using a diverse real-world dataset, our framework significantly improves navigation success rates, cover utilization, and threat avoidance compared to existing methods.